
\Problem{Base més petita amb el mateix nombre de dígits}

\Statement

Recordeu que en un sistema de numeració posicional, la base és el nombre de
dígits (símbols) diferents utilitzats per a representar els nombres. Per exemple, en
base 2, anomenat sistema de numeració binari, només existeixen dos dígits: 0 i
1. En base 10, el sistema de numeració decimal utilitzat habitualment,
existeixen 10 dígits (del 0 al 9). També es poden considerar bases superiors
a 10, com per exemple el sistema sexagesimal (emprat pels sumeris i pels
babilonis) que usa base 60.

\medskip

Donats dos nombres naturals, es vol trobar la base més petita (a partir de
dos) en la qual els dos nombres tenen el mateix nombre de dígits. Observeu que
aquesta base sempre existeix: agafeu el màxim, sumeu una unitat, i la
representació és d'un dígit en aquesta base.

\medskip

Per a fer-ho, escriviu i useu una acció

\begin{lstlisting}
        void base_amb_mateix_nombre_de_digits(int x, int y, int& b, int& n);
\end{lstlisting}

que, donats dos naturals $x$ i $y$,
desi en $b$ quina és la base més petita (amb $b\ge2$) en la qual $x$ i $y$ tenen el
mateix nombre de dígits, i desi en $n$ el nombre de dígits
de $x$ i $y$ en aquesta base.

\Input

L'entrada conté diferents casos.
Cada cas conté un parell de naturals, tots dos són menors o iguals que $2\cdot10^9$.

\Output

Per a cada cas de l'entrada, la base més petita en la qual tenen
el mateix nombre de dígits i el nombre de dígits que tenen en aquesta base.

\Hint

Fixeu-vos que quan $x$ i $y$ són molt allunyats l'un de l'altre, calcular
els dígits de l'un i de l'altre en base $b$ i veure si coïncideixen és lent.

\Observation

No podeu usar la biblioteca matemàtica per resoldre aquest problema.

\Sample
